% Feature gallery: one frame per technique a complex talk needs, each built
% so every variant (slides, notes, script, handout, trans, article) still
% comes out right. Copy the frame you need into your own talk.tex.
%
% Techniques, in order:
%   agenda from \tableofcontents          overlays: \pause, \item<n->, \only,
%   \alert<n>, \uncover, \visible           TikZ diagram revealed in steps
%   theorem / definition / example blocks  columns with an image
%   pgfplots with an overlay per series     booktabs table with a row reveal
%   code listing ([fragile] frame)          several notes per frame (\note[item])
%   mode-specific content (<handout:0>, \only<article>, \mode<article>)
%   a frame that exists only on slides      a button that jumps back to it
%   bibliography                            appendix with backup slides
\documentclass[
  beameroptions={aspectratio=169,ignorenonframetext,11pt},
  articleoptions={11pt,a4paper},
]{beamerswitch}
\usepackage{yjtalk}
\usepackage{listings}
\mode<presentation>{\usepackage{appendixnumberbeamer}}

\lstset{
  basicstyle=\ttfamily\footnotesize,
  keywordstyle=\color{yjWarm}\bfseries,
  commentstyle=\color{yjSecondary}\itshape,
  stringstyle=\color{yjAccent},
  columns=fullflexible,
  keepspaces=true,
  showstringspaces=false,
}

\title{Feature gallery}
\subtitle{Every technique a complex talk needs, from one source file}
\author{Teo Yu Jie}
\institute{beamerswitch template}
\date{29 September 2026}

\begin{document}

\begin{frame}[plain]
  \titlepage
  \note{This talk is a catalogue, not an argument: each frame shows one
    technique and says in its note how each output treats it.}
\end{frame}

\mode<article>{\maketitle\tableofcontents}

% --------------------------------------------------------------------------
\begin{frame}{Agenda}
  \tableofcontents
  \note{\texttt{\textbackslash tableofcontents} builds the agenda from
    \texttt{\textbackslash section} commands placed between frames.
    Sections also become PDF bookmarks and article headings.}
\end{frame}

\section{Overlays}

Overlays reveal a frame in steps. The slides and notes outputs keep every
step; the handout, transparencies and script collapse each frame to its final
state, so nothing is hidden on paper.

\begin{frame}{Overlays reveal a frame step by step}
  \begin{itemize}
    \item<1-> \texttt{\textbackslash item<2->} shows an item from step 2 on.
    \item<2-> \alert<2>{\texttt{\textbackslash alert<2>}} highlights on step 2 only.
    \item<3-> \only<3>{\texttt{\textbackslash only<3>} takes no space when hidden.}%
              \only<4>{\texttt{\textbackslash only<4>} replaces it on step 4.}
    \item<4-> \uncover<4->{\texttt{\textbackslash uncover} reserves space before it appears.}
  \end{itemize}
  \vspace{1ex}
  \visible<4>{\yjkey{Handout, trans and script show only the final step.}}
  \note{Click four times. On paper this frame is one page showing step 4,
    because handout mode collapses overlays.}
\end{frame}

\begin{frame}{\texttt{\textbackslash pause} steps through a derivation}
  \begin{align*}
    (a+b)^2 &= (a+b)(a+b) \\ \pause
            &= a^2 + ab + ba + b^2 \\ \pause
            &= a^2 + 2ab + b^2
  \end{align*}
  \pause
  \centering Each \texttt{\textbackslash pause} adds one step.
  \note{\texttt{\textbackslash pause} is the shortest overlay: everything after
    it appears one step later. Prefer \texttt{\textbackslash item<n->} in
    lists, since pauses do not nest.}
\end{frame}

\section{Diagrams and charts}

\begin{frame}{TikZ diagrams build up with the same overlay syntax}
  \centering
  \begin{tikzpicture}[
      box/.style={draw=yjBorder, fill=yjSurface, rounded corners, minimum width=2.6cm,
                  minimum height=1cm, font=\small},
      arr/.style={-Stealth, yjSecondary, thick}]
    \node[box] (src) {talk.tex};
    \node[box, right=1.4cm of src, visible on=<2->] (tex) {pdflatex $\times 7$};
    \node[box, right=1.4cm of tex, visible on=<3->, draw=yjAccent] (pdf) {seven PDFs};
    \draw[arr, visible on=<2->] (src) -- (tex);
    \draw[arr, visible on=<3->] (tex) -- (pdf);
    \node[below=0.4cm of tex, font=\footnotesize, text=yjSecondary, visible on=<2->]
      {jobname suffix picks the mode};
  \end{tikzpicture}
  \note{The \texttt{visible on=<n->} style is defined in yjtalk.sty; it works on
    any TikZ node or path. In the article the whole diagram shows at once.}
\end{frame}

\begin{frame}{Charts add one series per step}
  \centering
  \begin{tikzpicture}
    \begin{axis}[yj, width=11cm, height=5.5cm, xmin=0, xmax=4, ymin=0, ymax=16,
        xlabel={$x$}, ylabel={$y$}]
      \addplot[yjForeground, domain=0:4, samples=40] {x^2};
      \node[anchor=south east, font=\footnotesize] at (axis cs:3.3,10.9) {$x^2$};
      \only<2->{
        \addplot[yjWarm, domain=0:4, samples=40] {2^x};
        \node[anchor=north west, font=\footnotesize, text=yjWarm] at (axis cs:3.4,9.6) {$2^x$};
      }
      \only<3->{
        \fill[yjAccent] (axis cs:2,4) circle[radius=2.5pt];
        \node[anchor=north west, font=\footnotesize, text=yjAccentText] at (axis cs:2,4) {equal at $x=2$};
      }
    \end{axis}
  \end{tikzpicture}
  \yjsource{Source lines go in the footline via \texttt{\textbackslash yjsource\{...\}}.}
  \note{Wrap extra plots in \texttt{\textbackslash only<n->} inside the axis.
    For data from a CSV, see the Breeden--Litzenberger talk: it reads
    \texttt{data/curves.csv} with \texttt{\textbackslash pgfplotstableread}.}
\end{frame}

\section{Structured content}

\begin{frame}{Blocks mark definitions, theorems and examples}
  \begin{definition}[Convex]
    $f$ is convex if $f(tx + (1-t)y) \le t f(x) + (1-t) f(y)$ for $t\in[0,1]$.
  \end{definition}
  \begin{theorem}<2->[Jensen]
    For convex $f$, $f(\mathbb{E}X) \le \mathbb{E} f(X)$.
  \end{theorem}
  \begin{example}<3->
    $f(x)=x^2$ gives $(\mathbb{E}X)^2 \le \mathbb{E}X^2$, so variance is non-negative.
  \end{example}
  \note{Beamer's theorem environments also work in the article, where they
    become numbered theorems. Blocks take an overlay spec.}
\end{frame}

\begin{frame}{Columns put a figure beside its explanation}
  \begin{columns}[c]
    \column{0.45\textwidth}
      \includegraphics[width=\linewidth]{example-image}
    \column{0.5\textwidth}
      \begin{itemize}
        \item Keep images in \texttt{figures/} beside \texttt{talk.tex}.
        \item Prefer PDF or SVG-converted PDF over PNG.
        \item Let the chart fill its column; title carries the message.
      \end{itemize}
  \end{columns}
  \note{The image here is \texttt{example-image} from the mwe package.
    Replace it with \texttt{figures/your-file.pdf}.}
\end{frame}

\begin{frame}{Tables reveal one row at a time}
  \centering
  \begin{tabular}{lrr}
    \toprule
    Output  & Pages here & Overlays \\
    \midrule
    slides, dark & one per step  & kept \\ \pause
    notes   & one per step  & kept, notes on the right \\ \pause
    handout & 3 frames each & collapsed \\
    script  & one per frame & collapsed \\
    \bottomrule
  \end{tabular}
  \note{\texttt{\textbackslash pause} works between table rows. Keep
    \texttt{\textbackslash bottomrule} after the last pause so the rule is
    always drawn.}
\end{frame}

\begin{frame}[fragile]{Code needs a fragile frame}
\begin{lstlisting}[language=Python]
def build(talk):
    # one source, seven outputs
    for variant in ("slides", "dark", "notes", "script",
                    "handout", "trans", "article"):
        compile(talk, jobname=f"talk-{variant}")
\end{lstlisting}
  \note{Any frame with verbatim or listings needs \texttt{[fragile]}. Keep the
    listing flush left: beamer writes the frame to a temporary file.}
\end{frame}

\section{Notes and modes}

\begin{frame}{Notes can be a list of cues}
  \begin{itemize}
    \item Use \texttt{\textbackslash note[item]\{...\}} for bullet cues.
    \item Several \texttt{\textbackslash note} calls on one frame stack up.
  \end{itemize}
  \note[item]{Time: 30 seconds.}
  \note[item]{Point at the second bullet.}
  \note[item]{Transition: which content belongs in which output?}
\end{frame}

\begin{frame}{Some content belongs in one output only}
  \begin{itemize}
    \item<handout:0> This line is on the slides but not in the handout.
    \item \only<article>{Only readers of the article see this.}%
          \only<presentation>{\texttt{\textbackslash only<article>} text is hidden here.}
    \item Text between frames is article-only (\texttt{ignorenonframetext}).
  \end{itemize}
  \note{Mode names: beamer (slides, dark, notes), handout (handout, script), trans,
    article; presentation means all but article. Put \texttt{<handout:0>} on
    anything you do not want printed.}
\end{frame}

\mode<article>{
  \paragraph{Article-only block.} Wrapping content in
  \texttt{\textbackslash mode<article>\{...\}} is how you add a longer explanation, a proof, or
  a figure that would crowd a slide.
}

\begin{frame}<handout:0|trans:0>[label=live]{A frame that exists only on the live slides}
  \centering\Large Live demo goes here.
  \note{\texttt{\textbackslash begin\{frame\}<handout:0|trans:0>} drops this
    frame from the handout, the script and the transparencies: the script
    is built in handout mode.}
\end{frame}

\begin{frame}{References are ordinary \texttt{thebibliography}}
  \begin{thebibliography}{9}
    \bibitem{tantau} T.~Tantau, J.~Wright, V.~Miletić.
      \newblock \emph{The Beamer class}. User guide, CTAN.
    \bibitem{ball} A.~Ball.
      \newblock \emph{beamerswitch: convenient mode selection in Beamer documents}. CTAN.
  \end{thebibliography}
  \note{Cite with \texttt{\textbackslash cite\{tantau\}} anywhere. For long
    reference lists, put them in the appendix.}
\end{frame}

\begin{frame}{What to remember}
  \begin{enumerate}
    \item One \texttt{talk.tex}; the jobname picks the output.
    \item<2-> Overlays for the room, collapsed on paper.
    \item<3-> A note on every frame; the build refuses frames without one.
  \end{enumerate}
  \note{End here in a real talk. Everything after is appendix.}
\end{frame}

\appendix
\section*{Appendix}

\begin{frame}{Backup: jumping back to an earlier frame}
  Appendix frames are not counted in the total on the main slides.
  \par\medskip
  \hyperlink{live}{\beamerbutton{Back to the live demo}}
  \note{appendixnumberbeamer restarts numbering here, so the footer total
    matches the main talk. The button links to the frame labelled
    \texttt{live}; use it to answer a question without paging back.}
\end{frame}

\end{document}
